\documentclass[11pt]{article}
\usepackage{amsmath,hyperref}
\title{Critical rate for possible monitor outcomes under a common scalar shift}
\date{26 November 2026}
\begin{document}
\maketitle

\paragraph{Model and question.}
At strictly increasing cuts $t_i$ an automaton can switch states
according to an unknown projected shift $z_i$. For each transition
there is a finite union of bounded, nonempty scalar intervals on
which it is enabled; open and closed ends are distinguished.
Between cuts the law is
\[
 |z_{i+1}-z_i|\le B_{i+1}-B_i+\rho(t_{i+1}-t_i)
\]
for a known nondecreasing cumulative allowance $B_i$ and $\rho\ge0$.
This is a finite-cut uncertainty envelope, not a physical rate
identified from telemetry. Keep the correct automaton state when
changing the law. Ask for the least possible rate of reaching a
target state at the last cut, or its infimum if not attained.
When several evaluations actually share one constant shift, first
compile their joint input/output transitions using that common
value; closing interval ends separately before this compilation
could introduce nonexistent intermediate histories.

\paragraph{Fixed path criterion.}
Fix an accepting path and a connected enabling interval
$I_i$ with endpoints $\ell_i,u_i$ at each cut. For closed intervals,
a suitable sequence exists if and only if
$\ell_i-u_j\le D(i,j)$ for all pairs, where
\[
 D(i,j)=|B_i-B_j|+\rho |t_i-t_j|.
\]
The necessity follows by summing adjacent bounds. For sufficiency
take $z_i=\max_j\{\ell_j-D(j,i)\}$. It includes its own lower bound,
lies below the upper bound by the pair test, and satisfies adjacent
limits by the triangle inequality. The Lean file checks the finite
closed-interval equivalence, allowing a directed triangle bound;
this replay uses the symmetric instance. This is a scalar
Lipschitz selection argument, not a general selection theorem.
For context, see \href{https://arxiv.org/abs/1801.00325}{Fefferman
and Shvartsman, Sharp finiteness principles for Lipschitz selections}
for a much broader compact-convex setting. The new use here is
to invert a finite prefix-dependent transition law.

The original intervals satisfy the path criterion with one additional
rule: a cross-pair inequality must be strict if the selected lower
or upper endpoint is open. An equality would require both of those
specific endpoints. Conversely, shrink each open endpoint into its
interval by a common sufficiently small amount. Nonempty interval
lengths are preserved, and every cross-pair involving a shrunken
endpoint has positive slack. The closed criterion gives a witness
for the original intervals. A singleton is already closed when nonempty.

\paragraph{Finite inversion and exact endpoint test.}
For the closed relaxation of each path, the critical rate is the
maximum of zero and its signed endpoint ratios
\[
  \frac{\ell_i-u_j-|B_i-B_j|}{|t_i-t_j|},\qquad i\ne j .
\]
A finite set of possible paths and interval choices gives the
infimum over the closed relaxations. At every strictly larger rate
the optimal relaxed path's cross-pair bounds have positive slack,
so its original intervals also admit a run.
Compute all candidate ratios over labeled endpoints at different cuts
(together with zero), and find the first accepted relaxed candidate
by monotone search. Test the original intervals at this candidate
exactly to determine whether the critical rate is attained. If no
candidate accepts the relaxation, there is no accepting path at any
finite rate. The fixed-rate algorithm propagates unions of intervals
per state, expands them by the next allowance and intersects them
with the enabled bands, preserving endpoint flags. Each operation is
exact over declared parameter histories. The independent direct
implementation enumerates interval choices and evaluates both
the signed inequalities and their strictness. All 1{,}500 comparisons
agree.

\paragraph{Three-cut obstruction.}
Use the rectangular convention and latch law from the preceding
source replay: $x\in[0,1)$ at $y=1/2$, at least one return to hit,
initially active, and three misses needed for release. For
$z\in[0,5/2]$ and observed $x$ lists $\{-1,-2\}$, $\{-1\}$,
and $\{0,-1\}$ at times 0, 1 and 2, the miss sets are respectively
$[0,1)$, $[0,1)\cup[2,5/2]$, and $[2,5/2]$.
Every pair of misses fits rate $3/4$, but the three misses require
one adjacent jump from strictly below 1 to at least 2. Rate 1
is not enough; every strictly greater rate is. Equal endpoints in
the relaxed test would manufacture a release. In a suite switching
allowance at only one inter-cut edge, the same state relation must
be propagated; changing a rate is not a new episode.

The release modeled here is an output of the configured counter,
not confirmation of physical separation, independent measurements,
or authority to move. The replay does not measure or enforce a
deployed frame-change rate. No new contract disposition is claimed.
\end{document}
